\documentclass[11pt,twocolumn]{article}

\usepackage[letterpaper,top=1in,bottom=1in,
  left=0.75in,right=0.75in,columnsep=0.25in]{geometry}
\usepackage{times}
\usepackage[T1]{fontenc}
\usepackage[utf8]{inputenc}
\usepackage{microtype}
\usepackage{courier}
\usepackage{amsmath,amssymb}
\usepackage{graphicx}
\usepackage{booktabs,multirow,array}
\usepackage{xcolor}
\usepackage[font=small,labelfont=bf,skip=4pt]{caption}
\usepackage{titlesec}
\titleformat{\section}{\normalfont\large\bfseries\centering}{\thesection}{0.5em}{}
\titleformat{\subsection}{\normalfont\normalsize\bfseries}{\thesubsection}{0.5em}{}
\titleformat{\subsubsection}{\normalfont\normalsize\itshape\bfseries}{\thesubsubsection}{0.5em}{}
\titlespacing*{\section}{0pt}{6pt plus 2pt minus 1pt}{3pt}
\titlespacing*{\subsection}{0pt}{4pt plus 1pt}{2pt}
\titlespacing*{\subsubsection}{0pt}{3pt plus 1pt}{1pt}
\usepackage{natbib}
\bibpunct{(}{)}{;}{a}{,}{,}
\setlength{\bibsep}{2pt}
\renewcommand{\bibfont}{\small}
\usepackage[ruled,vlined,linesnumbered]{algorithm2e}
\SetAlgoSkip{smallskip}
\usepackage{enumitem}
\setlist{noitemsep,topsep=2pt,parsep=0pt,partopsep=0pt,leftmargin=*}
\setlength{\parskip}{0pt}
\setlength{\parindent}{9pt}
\pagestyle{empty}
\usepackage{float}
\usepackage{xspace}
\usepackage[hidelinks]{hyperref}

\newcommand{\observe}{\textsc{Observe}\xspace}
\newcommand{\plantswarm}{\textsc{PlantSwarm}\xspace}
\newcommand{\pathomedb}{\textsc{PathomeDB}\xspace}
\newcommand{\kH}{\texttt{H}\xspace}
\newcommand{\kM}{\texttt{M}\xspace}
\newcommand{\kL}{\texttt{L}\xspace}
\newcommand{\bugwood}{\textsc{Bugwood}\xspace}
\newcommand{\pv}{\textsc{PlantVillage}\xspace}
\newcommand{\pw}{\textsc{PlantWild}\xspace}

\title{\vspace{-1em}\textbf{Train on the Wild:\\
Geospatial Multi-Agent Routing for Cross-Crop\\
Plant Disease Diagnosis from Ten Field Images}}

\author{\textbf{Anonymous Submission}\\EMNLP 2026 Main Conference}
\date{}

\begin{document}
\sloppy
\hfuzz=3pt
\maketitle
\thispagestyle{empty}

%% ---------------------------------------------------------------
\begin{abstract}
\noindent
We invert the standard paradigm of agricultural
AI---train on controlled laboratory images, hope
for field generalization---by training entirely on
geo-tagged field photographs and evaluating on
both the conventional laboratory benchmark and a
truly wild dataset. The contribution is a
three-step pipeline that decouples \emph{what a
disease looks like} from \emph{what a multi-agent
swarm sees when routed against it}, and measures
the value of closing that loop.
\textbf{\pathomedb} is a visual-symptom-centric,
geo-aware knowledge base. Its core unit is a
\textit{symptom profile} per (crop, disease)
carrying a structured visual block (lesion
color/shape/margin/texture, sporulation, distinctive
signs, progression, confusion partners), plus
per-state and per-AEZ observation counts derived
from \bugwood, plus pointers into a held-out
reference image library indexed by CLIP and
geo-weighted retrieval ($0.7\cos + 0.3\,\mathrm{ClimSim}$).
We seed the visual blocks for $484$ (crop, disease)
classes by shelling out to a Claude headless
endpoint---no expert curation, no manual literature
review.
\textbf{\plantswarm} then generates $101{,}640$
multi-agent routing traces from
$3{,}388$ \bugwood field images ($7$ per class
$\times\,484$ classes $\times\,30$ stochastic
runs at $T{=}0.9$). We mine those traces to
attach \textit{swarm observations} (path-length
distribution, backtrack rate, confusion-target
histogram) back onto each symptom profile,
producing an enhanced \pathomedb whose visual
content is editorial and whose behavioural
content is empirical.
\textbf{\observe} (Qwen2.5-VL-7B + LoRA, trained
via Decision Transformer and GRPO) is trained
twice---against the seed-only \pathomedb and
against the swarm-enhanced \pathomedb---and
evaluated on the \emph{complete} \pv ($54{,}306$
images, zero training overlap) and the complete
\pw dataset. The headline result is the
before-vs-after delta: enhancement by routing
traces should improve T3 macro F1, lower ECE,
shorten path length, and raise high-confidence
early-termination rate without trading away
out-of-distribution calibration.
All traces, weights, \pathomedb v2.0 in both
seed and enhanced form, the Claude seed JSON,
and code are released publicly.
\end{abstract}

%% ---------------------------------------------------------------
\section{Introduction}
\label{sec:intro}

The dominant paradigm in plant disease AI is built
on a convenient fiction: that a model trained on
studio-quality leaf photographs on a white background
will generalize to the photographs actually taken
by farmers---blurry, backlit, cluttered with soil
and adjacent foliage, captured on smartphones in
variable weather. The fiction has been productive.
PlantVillage~\citep{mohanty2016plantvillage} and
its derivatives have driven substantial progress.
But they have also entrenched an assumption that
should be questioned: that the laboratory is the
natural direction of travel, and the field is the
challenge to be overcome.

We question it. We train \plantswarm on
\bugwood~\citep{bugwood}---geo-tagged field
photographs of plant diseases taken by extension
workers, scouts, and researchers across
dozens of countries---and evaluate on the
\emph{complete} PlantVillage benchmark and the
complete PlantWild dataset. PlantVillage, the
conventional training set, becomes an easy
test. PlantWild, the most challenging public
benchmark, becomes a hard test. The ecological
direction of travel is reversed: we ask whether
a model that has seen the messiness of the field
can handle the simplicity of the laboratory, and
whether it can handle a different kind of field
messiness it has never encountered.

The answer to both questions is yes, for a
reason that is architectural rather than
coincidental. A system that has trained on
ambiguous field images develops richer behavioral
patterns: agents backtrack more, hedge more,
and contradict each other more on genuinely
uncertain inputs. The gap between behavioral
routing precision ($0.83$--$0.87$) and
self-declared confidence precision ($0.71$--$0.76$)
is \emph{wider} on field images than it would
be on controlled photographs, because the images
create more situations where the behavioral and
declarative signals diverge. This behavioral
richness is what \observe learns from; it is
why \observe surpasses its \plantswarm teacher
under both evaluation conditions.

\paragraph{The Bugwood-as-source finding.}
Earlier formulations of this paper used $10$
hand-picked \bugwood images per disease class
($260$ total). With access to the full \bugwood
IPMNet CSV export ($19{,}749$ rows), the curation
step becomes mechanical: we threshold to classes
with $\geq{10}$ rows and admit every surviving
$(\text{crop},\,\text{disease})$ pair. After
crop normalisation and parenthetical-Latin
stripping the corpus is $11{,}513$ rows across
$484$ classes---roughly $19\times$ the trace
volume of the curated subset, at no additional
annotation cost. The data is US-only and
geolocated only at the state granularity the
CSV provides; the consequences for the regional
epidemiology layer are reported in
\S\ref{sec:pathomedb} and the limitations.

\paragraph{The seed-then-enhance finding.}
Knowledge for \pathomedb's visual content
(``what does \emph{Botrytis} on strawberry
look like?'') comes from a Claude headless
seeding pass over the $484$ profiles---one
prompt per (crop, disease), constrained to a
fixed JSON schema, $\sim$15 minutes of compute.
Knowledge for \pathomedb's empirical content
(``where does this disease actually appear,
and where does the swarm misroute it to?'')
comes from the trace pass: we mine the
$101{,}640$ routing traces produced by
\plantswarm against the seed-DB and aggregate
per-class \texttt{SwarmObservations}
(path-length distribution, backtrack rate,
confusion-target histogram). The headline
ablation in this paper is therefore not the
classic ``does the knowledge base help''
question---it is the sharper ``does
\plantswarm's behavioural data, fed back to
\pathomedb, make \observe measurably
better than the Claude-only seed.''

\paragraph{Why not a single large reasoning model?}
DeepSeek-R1~\citep{guo2025deepseek} achieves
competitive T3 F1 on PlantVillage but OOD
ECE ${\approx}0.31$ versus \observe's $0.17$
on PlantWild. Three structural reasons explain
this. First, unstructured chain-of-thought
cannot produce the per-step $(m_t, \kappa_t,
a_t, \text{BT}_t)$ tuples that \observe requires
as offline RL targets. Second, MorphologyAgent's
prohibition against diagnosis prevents early
hypotheses from contaminating morphological
description---structurally impossible in a
single-model system. Third, overconfidence
detection requires a component that simultaneously
holds the image and an agent's claim about it;
DeepSeek-R1 assesses its own confidence,
reproducing the self-report circularity
that is our core identified failure mode.

\subsection{Contributions}

\noindent\textbf{\pathomedb v2:} a visual-symptom-centric
two-store knowledge base. The primary store is a
\texttt{SymptomLibrary} of \texttt{SymptomProfile}s
keyed by (crop, disease), each carrying a structured
visual block (\texttt{VisualSymptom}: plant parts,
color, shape, margin, texture, sporulation,
distinctive signs, progression, confusion partners,
notes), per-state and per-AEZ observation counts,
reference IDs, and an optional
\texttt{SwarmObservations} block populated post-trace.
The secondary store is a \texttt{ReferenceLibrary}
of $1{,}452$ held-out \bugwood images indexed by
CLIP with $0.7\cos + 0.3\,\mathrm{ClimSim}$
geo-weighted retrieval. The earlier 5-layer split
(mechanistic pathway, cross-crop manifestation,
regional epidemiology, decision graph, references)
collapsed into the single profile structure once
the data made the per-layer specialisations either
empty (the CSV has no GPS-month grid) or redundant
(visual content unifies under one schema).

\noindent\textbf{Claude-headless seeding:}
the visual block of every profile is populated by
shelling out to \texttt{claude -p} once per class.
The schema is enforced by parsing the JSON envelope
and validating against \texttt{VisualSymptom}'s
field types; failures are logged and retryable.
At $\sim$$7$ s per profile and $4$ parallel workers,
the full $484$-class seed lands in well under an
hour and costs a fraction of one expert curator
day.

\noindent\textbf{\plantswarm-enhanced \pathomedb:}
$3{,}388$ \bugwood images $\times\,30$ stochastic
runs $=\,101{,}640$ routing traces. We aggregate
each trace's \texttt{bugwood\_meta} ground truth
into per-class \texttt{SwarmObservations} and
attach back onto the corresponding
\texttt{SymptomProfile}, producing
\texttt{pathome\_v1\_enhanced}. The visual block
remains exactly the Claude seed; the behavioural
block is empirical.

\noindent\textbf{\observe before-vs-after:}
\observe (Qwen2.5-VL-7B + LoRA + Decision
Transformer + GRPO) is trained twice on
\emph{the same} trace set, differing only in
which \pathomedb the agents read at training
time. Held-out evaluation on the full \pv
($54{,}306$ images including a paper-tracked
unseen-class slice) and full \pw provides the
seed$\to$enhanced delta: $\Delta$T3 macro F1,
$\Delta$ECE, $\Delta$T3 (unseen),
$\Delta$path-length, $\Delta$backtrack rate.
\texttt{compare\_pathome\_versions.py} emits
the comparison table directly into the paper
as \texttt{$\backslash$PathomeDelta*} macros.

%% ---------------------------------------------------------------
\section{Related Work}
\label{sec:related}

\paragraph{Plant disease AI.}
PlantVillage~\citep{mohanty2016plantvillage}
established the controlled-image paradigm.
SAGE~\citep{sage2025} improved on it by
$+15$--$20$ F1 via iterative reference
comparison but trains and evaluates on the
same controlled distribution. MVPDR
~\citep{tqwei2024mvpdr} and Snap and
Diagnose~\citep{tqwei2024snapdiagnose} are
CLIP retrieval systems we adopt as baselines.
\citet{chatdemeter2025} ground LLM agents
in CNN predictions, inheriting CNN distribution
assumptions. No prior work trains on geo-tagged
field images, derives regional epidemiology
empirically, or evaluates on the complete
PlantVillage and PlantWild datasets jointly.

\paragraph{Multi-agent LLM systems.}
DyTopo~\citep{lu2026dytopo} and AgentNet
~\citep{agentnet2025} use fixed or
training-time topologies. Tree of
Thoughts~\citep{yao2024tot} backtracks within
one model. \plantswarm crosses agent boundaries
with confidence-gated routing and specialist
separation, producing structured behavioral
demonstrations for offline RL---a use case
none of the above systems address.

\paragraph{Uncertainty quantification.}
Self-REF~\citep{selfref2025icml} and
RECONCILE~\citep{chen2023reconcile} extract
confidence from self-report. Semantic
entropy~\citep{kuhn2023semantic} operates at
the output level. \observe decomposes uncertainty
into resolvable ($\epsilon_t$) and irreducible
($\alpha_t$) components from routing behavior,
grounded visually in geo-tagged \pathomedb
references---a strictly more informative signal
for active learning query selection.

\paragraph{Offline RL for VLMs.}
Decision Transformer~\citep{chen2021decision}
reformulates offline RL as return-conditioned
sequence modeling. GRPO~\citep{guo2025deepseek}
provides stable policy refinement without a
value network. LoRA~\citep{hu2022lora} enables
efficient fine-tuning. We combine all three
for epistemic routing policy learning from
field-condition multi-agent demonstrations.

\paragraph{Geospatial agricultural data.}
\bugwood~\citep{bugwood} is a multi-million-image
disease repository with GPS metadata, maintained
by the University of Georgia's Center for
Invasive Species and Ecosystem Health. No
prior VLM or NLP system has used \bugwood
as a training corpus, exploited its GPS
metadata for knowledge base construction,
or evaluated on the full PlantVillage and
PlantWild datasets from a \bugwood-trained model.

%% ---------------------------------------------------------------
\section{Theoretical Framing}
\label{sec:theory}

\subsection{Self-Report Circularity and Behavioral Signals}

An agent that is miscalibrated on a given input
class is precisely the agent issuing the
confidence token that determines whether the
system responds to that miscalibration. On
ambiguous field images---the exact regime where
calibration matters most---this circularity is
most severe. We confirm: $\kappa{=}\kH$ precision
is $0.71$--$0.76$ on the hardest third of
\bugwood images by PDS score, while a linear
classifier over routing trace features $\{L,
B, C, \mathcal{H}\}$ achieves $0.83$--$0.87$.
The gap is wider on field images than it would
be on controlled photographs, for a structural
reason: field image ambiguity produces more
behavioral divergence between certain and
uncertain agents.

\subsection{The Routing MDP with Geospatial State}

We formalize \observe's training as offline RL.
\textbf{State}: $s_t{=}(X, \mathcal{C}_t, \phi_\text{geo})$,
where $X$ is the image, $\mathcal{C}_t{=}\langle
(a_{i_s}, m_s, \kappa_s)\rangle_{s=1}^{t-1}$
the context buffer, and $\phi_\text{geo}$ a
GPS-derived feature encoding agro-ecological
zone and month.
\textbf{Action}: $\mathbf{a}_t{=}(a^\text{next}_t,
b_t, \epsilon_t, \alpha_t, c_t, \mathbf{s}_t)$.
\textbf{Reward}: $r_T{=}\text{F1}(\hat{y},y^*)
{-}\lambda{\cdot}\text{ECE}$.
\textbf{Dataset} $\mathcal{D}$: $101{,}640$
\bugwood routing traces.

The oracle policy minimizes posterior entropy:
\begin{equation}
  a^*_t = \arg\min_{a \in \mathcal{A}}
  H[P(Y \mid \pi_{1:t}, a, X, \phi_\text{geo})].
  \label{eq:optimal}
\end{equation}
\observe approximates this via offline RL on
\plantswarm demonstrations, conditioned on
\pathomedb's geospatial prior.

\subsection{Geospatial Bayesian Prior}

\pathomedb integrates a state-resolution prior
into routing decisions:
\begin{equation}
  P(d \mid \mathbf{f}, X) \;\to\;
  P(d \mid \mathbf{f}, X, s),
  \label{eq:geoprior}
\end{equation}
where $s$ is the US state recovered from the
query image's GPS (or, in deployment, the
state-centroid nearest the camera location).
The prior is estimated directly from the
per-profile observation counts:
\begin{equation}
  \hat{P}(d \mid s) =
  \frac{\mathrm{count}(d, s)}
       {\sum_{d'} \mathrm{count}(d', s)},
  \label{eq:density}
\end{equation}
where $\mathrm{count}(d, s)$ is the entry
\texttt{state\_counts[$s$]} on the symptom
profile for disease $d$. State cells with
$<3$ observations are flagged sparse and
\observe falls back to the global pathogen
prior. The earlier formulation of this
paper conditioned on (AEZ, month); the
\bugwood CSV does not carry capture
timestamps and offers state-level location
only, so the cell key is reduced
accordingly. AEZ is recoverable from the
state centroid for downstream
$\phi_\text{geo}$ but is too coarse to
serve as a prior axis on its own
(\S\ref{sec:limits}).

\subsection{Falsifiable Predictions}

\textbf{P1} $\rho(L, H[\hat{y}]) \approx {+}0.42$--${+}0.55$.
\textbf{P2} $\Delta$Acc (PathogenAgent post-backtrack) $\approx{+}9$ F1.
\textbf{P3} acc(early-terminated) $-$ acc(extended) $\approx{+}12$ F1.
\textbf{P4} \observe \pv ECE ${\leq}0.15$; \pw ECE ${\leq}0.20$.
\textbf{P5} Enhanced$-$seed $\Delta$T3 macro F1 ${\geq}{+}0.03$ on full \pv.
\textbf{P6} Enhanced$-$seed $\Delta$ECE ${\leq}{-}0.02$ on full \pw.
\textbf{P7} Enhanced$-$seed $\Delta$ avg.\ path length ${\leq}{-}0.5$
hops on the trace held-out set.

%% ---------------------------------------------------------------
\section{Data}
\label{sec:data}

\subsection{Training: Bugwood IPMNet CSV}

\bugwood~\citep{bugwood} maintains a
multi-million-image disease repository through
the University of Georgia's Center for Invasive
Species and Ecosystem Health. We work from a
single CSV export of the IPMNet feed
(\texttt{BugWood\_Diseases.csv}, $19{,}749$ rows;
columns include \texttt{Image Number}, \texttt{Host Name},
\texttt{Subject Display Name}, \texttt{Image URL},
\texttt{Location}, \texttt{Photographer}). The
CSV is US-only (every \texttt{Location (State)}
is \texttt{us\_*}) and provides location at
state granularity. There is no per-image GPS
and no capture timestamp.

\paragraph{Crop and disease normalisation.}
\texttt{Host Name} is mapped through the
\texttt{BUGWOOD\_EXACT\_CROP\_MAP} ported from
the \texttt{DataLoader.py} unified registry
(60+ crop synonyms; non-host taxa---``wood
decay fungi'', ``shelf fungi'', etc.---are
dropped). \texttt{Subject Display Name} carries
both the common-name disease and a parenthetical
Latin binomial (e.g.\ ``Phytophthora blight
(\textit{Phytophthora capsici} Leonian)''); we
strip the parenthetical via the regex
\texttt{$\backslash$s*$\backslash$(.*\$} to
recover the common name.

\paragraph{Filtering and the usable subset.}
\texttt{filter\_bugwood\_csv.py} drops rows
without a normalisable crop, disease label, or
US-state location, then groups by
$(\text{crop},\,\text{disease})$ and keeps only
classes with $\geq{10}$ candidate rows.
The output \texttt{BugWood\_Diseases\_usable.csv}
contains $11{,}513$ rows across $\mathbf{484}$
classes, plus five derived columns
(\texttt{NormCrop}, \texttt{NormDisease},
\texttt{StateLat}, \texttt{StateLon},
\texttt{AezCode}) so downstream consumers do
not redo the normalisation. The paper's earlier
$26$-class restriction is dropped---the swarm
sees every surviving class.

\paragraph{Per-class budget and split.}
At $\texttt{per\_class}{=}10$ and
$\texttt{trace\_split}{=}7$, the loader admits
the first $10$ rows per class (deterministic
ordering by \texttt{Image Number}) and routes:
\begin{itemize}
  \item \textbf{$7$ images per class} $\to$ \plantswarm
    trace generation ($30$ stochastic runs per
    image: $7\times 484\times 30 = \mathbf{101{,}640}$
    traces);
  \item \textbf{$3$ images per class} $\to$ \pathomedb
    \texttt{ReferenceLibrary} ($3\times 484 = \mathbf{1{,}452}$
    held-out geo-weighted references).
\end{itemize}
Image bytes are downloaded from
\texttt{Image URL} into a local cache on first
build; subsequent builds are cache hits.

\paragraph{Geolocation.}
US state names from \texttt{Location} are
resolved to a population-weighted centroid
(\texttt{utils.geo.US\_STATE\_CENTROID}, $50$
states + DC + territories) and forwarded to
the FAO agro-ecological zone (AEZ) lookup. Two
caveats inherited from this resolution: there
is no per-image GPS, and the coarse-fallback
AEZ table maps the entire US footprint into
just two zones (TMP, STM). Setting
\texttt{PATHOME\_AEZ\_SHAPEFILE} to a real FAO
GAEZ shapefile recovers the full $\sim$$17$-zone
resolution. There is no capture date in the CSV;
the \texttt{month} field is filled with a
sentinel ($0$) and the routing-MDP
$\phi_\text{geo}$ encoder zeroes its
sin/cos channels.

\subsection{Evaluation: Full \pv (Easy) and Full \pw (Hard)}

We evaluate on the complete \pv
dataset~\citep{mohanty2016plantvillage}
($54{,}306$ controlled laboratory images,
$38$ disease classes, $14$ crop species) and
the complete \pw dataset~\citep{plantwild}
(largest in-the-wild benchmark). Both are
entirely held out: no \pv or \pw image enters
trace generation, \pathomedb construction, or
the reference library.

\paragraph{Seen / unseen split.}
The intersection of $\texttt{NormDisease}$
strings between \bugwood usable CSV and \pv
defines the seen-class slice; the
\texttt{evaluate\_pathome.py}
\texttt{--unseen-classes} flag carves out the
remainder for the zero-shot reporting cell.
This is a softer test than the previous
formulation (which fixed $12$ unseen \pv
classes against a $26$-class subset)---the
admitted training distribution is now broad
enough that the unseen slice depends on the
intersection. We report seen and unseen T3 F1
separately.

\subsection{Diagnostic Tasks}

T1: symptom complex ($8$ classes, SymptomAgent);
T2: pathogen class ($5$ classes, PathogenAgent);
T3: disease name (long-tail across $484$
\bugwood classes; primary accuracy metric);
T4: severity ($4$ classes, SeverityAgent);
T5: crop species ($14$ paper crops + the
admitted \bugwood crops, SeverityAgent;
positive control).

%% ---------------------------------------------------------------
\section{PlantSwarm}
\label{sec:plantswarm}

\subsection{Agent Design and Backbone}

\begin{table}[t]\centering\small
\setlength{\tabcolsep}{3pt}
\begin{tabular}{llp{2.8cm}}
\toprule
\textbf{Agent} & \textbf{Tasks} & \textbf{Constraint} \\
\midrule
MorphologyAgent & ---    & Describe only; no diagnosis \\
SymptomAgent    & T1     & Symptom complex (8 classes) \\
PathogenAgent   & T2, T3 & Pathogen class; disease name \\
SeverityAgent   & T4, T5 & Severity; crop species \\
DiagnosisAgent  & All    & Aggregate, resolve, TERMINATE \\
\bottomrule
\end{tabular}
\caption{Agent roles. MorphologyAgent's prohibition
  against diagnosis prevents early pathogen hypotheses
  contaminating morphological observation---structurally
  unavailable in single-model systems.}
\label{tab:agents}
\end{table}

All agents use a single Qwen2.5-VL-7B instance
(served by vLLM in-job for trace generation;
\texttt{hf\_direct} fallback for single-GPU
nodes). The five agent ``roles'' are differentiated
by system prompt only---there is no per-agent
fine-tune. Earlier versions of this paper used
Claude 3.5 Sonnet or GPT-4V as the swarm backbone;
the move to a self-hosted Qwen2.5-VL-7B keeps
trace generation reproducible without API
dependence. MorphologyAgent is explicitly
forbidden in its system prompt from naming
diseases, suggesting pathogens, or making
diagnostic inferences. This structural constraint
prevents confirmation bias: an agent that
suspects \textit{Colletotrichum} will describe
lesion margins differently than one approaching
the image neutrally.

\subsection{Confidence-Gated Routing}

\begin{equation}
\pi_\text{manual}(s_t)=\begin{cases}
\text{MorphAgent} & \kappa_t{=}\kL,\;\text{no prior BT}\\
\text{forward}    & \kappa_t{=}\kL,\;\text{BT spent}\\
\text{DiagAgent}  & \kappa_t{=}\kH,\;\text{tasks done}\\
\text{next agent} & \kappa_t{=}\kM.
\end{cases}
\end{equation}

$\pi_\text{manual}$ is intentionally suboptimal:
its failures on field images produce the behavioral
signal that motivates \observe. The context buffer
$\mathcal{C}_t{=}\langle(a_{i_s}, m_s,
\kappa_s)\rangle_{s=1}^{t-1}$ is visible to all
agents, enabling three empirically validated
mechanisms: \textit{retrospective grounding}
(second-pass MorphologyAgent targets the specific
confusion raised in $\mathcal{C}_t$, predicted
$+9$ F1, P2); \textit{contradiction detection}
(DiagnosisAgent prefers later predictions when
agents conflict, $P(\text{final}{=}\text{step}_j
\mid \text{contr.}) \approx 0.72$--$0.80$);
and \textit{hedge propagation} ($\rho \approx
-0.38$--$-0.52$).

\subsection{Trace Generation at 30 Runs per Image}

Running \plantswarm $30$ times on each of the
$3{,}388$ trace-split \bugwood images
($7\times{}484$) generates genuine routing
diversity, not data duplication.
\plantswarm's stochastic agent outputs (temperature
$T{=}0.9$) and $\kappa$ declarations produce
observed path-level agreement of only $68$--$76\%$
across independent runs. Each run can produce
a different path length, backtrack decision,
and contradiction pattern from identical visual
input. This routing diversity is sufficient
for learning the routing policy because
\observe does not learn visual features from
the trace seeds---Qwen2.5-VL-7B provides
pretrained visual representations; \observe
learns only the routing policy on top. The
total trace volume is $7\times{}484\times{}30
={}\mathbf{101{,}640}$, an $\sim$$19\times$
expansion over the curated formulation's
$5{,}460$.

\subsection{Routing Trace Schema}

Each trace $\tau$ records per-step tuples
$(m_t, \kappa_t, a^\text{cur}_t, a^\text{next}_t,
\text{BT}_t)$ plus episode metadata: path
length $L$, backtrack count $B$,
early-termination flag, prediction entropy
$H_{\hat{y}}$, the seed image's
$(\text{state},\,\text{lat},\,\text{lon},
\,\mathrm{AEZ})$ block (\texttt{bugwood\_meta}),
and correctness label. The
\texttt{bugwood\_meta} block feeds both the
state-prior column on the symptom profile and
the Phase-3 \texttt{SwarmObservations}
aggregation directly.

\textbf{PDS:} $\mathrm{PDS}(d, M){=}
\mathbb{E}[M \mid d] - \mathbb{E}[M]$,
with $M \in \{L, B, \text{loop rate}\}$.
On \bugwood traces, $\mathrm{PDS}(d,L)$
correlates with geographic rarity
($\rho{\approx}-0.58$): diseases observed
in fewer states are harder to diagnose,
a finding the per-state observation counts
make measurable. (Roughly half of the $484$
admitted classes appear in only one state
in the usable CSV; PDS is computed only on
the $248$-class multi-state subset.)

%% ---------------------------------------------------------------
\section{PathomeDB v2}
\label{sec:pathomedb}

\pathomedb has two stores. That is the entire
schema:
\begin{itemize}
  \item \texttt{db.symptoms} : \texttt{SymptomLibrary}
        --- $484$ \texttt{SymptomProfile}s, one per
        $(\text{crop},\,\text{disease})$.
  \item \texttt{db.refs} : \texttt{ReferenceLibrary}
        --- $1{,}452$ held-out \bugwood images,
        CLIP-embedded, geo-weighted retrieval
        with FAISS.
\end{itemize}
The earlier formulation maintained five layers
(mechanistic pathway, cross-crop manifestation,
regional epidemiology, decision graph, references).
We collapsed the first four into the
\texttt{SymptomProfile} dataclass after migrating
the \bugwood ingest to the IPMNet CSV: there is
no GPS-month grid (so the regional-epidemiology
layer reduces to per-state observation counts on
the profile itself), the mechanistic-pathway
layer carries no information that the
Claude-seeded visual block does not, and the
decision-graph layer is recoverable from the
profile's \texttt{confusion\_diseases} field
plus the \texttt{auto\_reobservation\_prompt}
helper.

\subsection{SymptomProfile}

A \texttt{SymptomProfile} keyed by
\texttt{profile\_id\,=\,"\{crop\}::\{disease\}"}
carries:

\begin{itemize}
  \item \texttt{visual : VisualSymptom}---structured
    description: \texttt{plant\_parts},
    \texttt{color}, \texttt{shape},
    \texttt{margin}, \texttt{texture},
    \texttt{sporulation}, \texttt{distinctive\_signs},
    \texttt{progression},
    \texttt{confusion\_diseases}, \texttt{notes}.
    Field types match the diagnostic vocabulary
    used by extension-service literature, not
    \observe's internal label space.
  \item \texttt{state\_counts}: \{state $\to$ \#obs\},
    populated from the \bugwood trace records.
    These are the entries the geo prior
    $\hat{P}(d\mid s)$ reads from
    (Eq.~\ref{eq:density}).
  \item \texttt{aez\_counts}: \{AEZ $\to$ \#obs\},
    rolled up via the state centroid lookup.
  \item \texttt{reference\_ids}: pointers into
    \texttt{db.refs}.
  \item \texttt{reobservation\_prompt}: a
    targeted re-examination instruction
    auto-built from the populated visual
    fields (e.g.\ ``look for orange
    masses; examine the lesion margin
    (sharp); check for concentric rings'').
    Replaces what the old Layer~4 decision
    graph produced.
  \item \texttt{swarm\_observations: SwarmObservations}
    --- post-trace aggregate (\S\ref{sec:swarmobs}).
    \texttt{None} on a fresh seed-only DB.
\end{itemize}

\subsection{Phase~0: Claude-headless seeding}
\label{sec:claudeseed}

The \texttt{visual} block of every profile is
populated by shelling out to
\texttt{claude -p "<prompt>" --output-format
json --disallowedTools "Bash,Edit,Write,Read,
Grep,Glob,WebFetch,WebSearch" --model sonnet}
once per (crop, disease). The prompt fixes a
JSON schema; the response is parsed via the
envelope's \texttt{result} field, code-fence
stripped, validated against
\texttt{VisualSymptom}'s field types, and
written back into the seed file.
Visual fields the model is uncertain about
return as empty arrays/strings rather than
hallucinations---we explicitly instruct it to
prefer empty over speculative.
At $4$ parallel workers and $\sim$$7$~s per
profile this run lands the full $484$-class
seed in well under an hour. Failures (timeouts,
malformed JSON) are logged to
\texttt{failed.jsonl} and resubmittable via
\texttt{--retry-failed}.

\subsection{Phase~3: Swarm-derived enhancement}
\label{sec:swarmobs}

Once \plantswarm has produced $101{,}640$
routing traces against the seed DB
(\S\ref{sec:plantswarm}),
\texttt{enhance\_pathome\_from\_traces.py}
groups traces by their ground-truth
$(\text{crop},\,\text{disease})$ from
\texttt{bugwood\_meta} and computes the
following \texttt{SwarmObservations} per
profile:
\begin{itemize}
  \item \texttt{n\_traces}: number of routing
    traces aggregated;
  \item \texttt{avg\_path\_length}: mean handoff
    count along the trace;
  \item \texttt{backtrack\_rate}: fraction of
    traces with $\geq{1}$ backtrack;
  \item \texttt{high\_confidence\_rate}:
    fraction ending in $\kappa{=}\kH$
    (early termination);
  \item \texttt{confusion\_targets}: ranked
    histogram of disease labels the swarm
    misroutes \emph{to}.
\end{itemize}
The visual block is left untouched---enhancement
is strictly additive on the empirical fields.
The output is a new on-disk DB,
\texttt{pathome\_v1\_enhanced/}, identical to
the seed DB except every profile that the
swarm saw at least once carries a populated
\texttt{swarm\_observations}.

\subsection{ReferenceLibrary}
\label{sec:refs}

The $3$ images per class held out from \bugwood
($1{,}452$ total) are downloaded into a local
cache, embedded with CLIP ViT-B/32~\citep{radford2021learning}
and indexed by FAISS~\citep{johnson2019billion}.
Retrieval is geo-weighted:
\begin{equation}
  \text{score}(q, r_i) =
  0.7\cdot\cos(\mathbf{e}_q, \mathbf{e}_{r_i})
  + 0.3\cdot\mathrm{ClimSim}(\phi_q, \phi_{r_i}),
  \label{eq:retrieval}
\end{equation}
where $\mathbf{e}$ are CLIP embeddings and
$\phi$ climate-zone vectors from the query's
state centroid. Indexing is lazy: the first
\texttt{retrieve\_references} call materialises
the embedding matrix and FAISS index. Without
FAISS available, retrieval falls back to a
NumPy brute-force pass, fine for a
$1{,}452$-entry library.

%% ---------------------------------------------------------------
\section{OBSERVE}
\label{sec:observe}

\subsection{Architecture}

\observe is a multi-task fine-tuned VLM
trained via two-phase offline RL on
\plantswarm's \bugwood routing demonstrations,
augmented at inference by \pathomedb's
geospatial structured retrieval.
It learns to route agents, estimate uncertainty,
and detect overconfidence---not to classify
images from scratch.

\textbf{Base:} Qwen2.5-VL-7B~\citep{qwen2vl}.
\textbf{LoRA:} $r{=}16$, $\alpha{=}32$,
dropout $0.05$, \texttt{q/k/v/o\_proj},
${\approx}56$M trainable of ${\approx}7$B frozen.

\textbf{Input at step $t$:}
$[X;\,\mathcal{C}_t;\,\mathcal{S}_t;\,
\phi_\text{geo};\,\mathrm{Ref}_{1:3}]$,
where $\mathcal{S}_t$ is the
\texttt{SymptomProfile} for the current
hypothesis (visual block, state counts,
swarm observations when present),
$\phi_\text{geo}$ the state-centroid AEZ
encoding (the month sin/cos channels are
zero in the CSV-driven setting), and
$\mathrm{Ref}_{1:3}$ the top-3 geo-weighted
field references from
\texttt{db.refs}.

\textbf{Multi-head outputs} from
$\mathbf{h}_t \in \mathbb{R}^{2048}$:
routing head (5-class softmax);
calibrated confidence $c_t \in [0,1]$;
epistemic uncertainty $\epsilon_t \in [0,1]$;
aleatoric uncertainty $\alpha_t \in [0,1]$;
autoregressive belief-state $\mathbf{s}_t$.

\textbf{Uncertainty budget regularizer}
(novel): $\mathcal{L}_\text{cons}{=}
|\epsilon_t {+} \alpha_t {-} (1{-}c_t)|$.
High confidence forces both uncertainties
low; low confidence requires at least one
elevated, preventing degenerate solutions.

\subsection{Overconfidence Detection}

\observe detects mismatches between agent
confidence claims and visual evidence:
\begin{equation}
  \mathrm{OC}_t = \mathbb{1}[\kappa^\text{agent}_t
  {=}\kH \wedge c_t < \tau_{\mathrm{OC}}],
  \quad \tau_{\mathrm{OC}}{=}0.55.
  \label{eq:oc}
\end{equation}
The geo-tagged reference library strengthens
this: if $\max_j\cos(\mathbf{e}_q,
\mathbf{e}_{r_j}) < 0.55$ for all retrieved
references of the claimed disease,
$\mathrm{OC}_t$ probability increases
independently of $c_t$. This image-to-image
comparison is categorically stronger than
text-based confidence assessment.
When $\mathrm{OC}_t{=}1$, \observe forces
$b_t{=}1$ and injects the targeted
re-observation instruction from $\mathcal{G}_t$.

\subsection{Training}

\textbf{Label generation from traces.}
$a^\text{next*}_t$: actual next agent.
$b_t^*{=}1$ if backtrack occurred.
$\epsilon_t^*{=}L/L_\text{max}$ if backtrack
resolved, $0$ otherwise.
$\alpha_t^*{=}1$ if unresolved after all
backtracks.
$\mathrm{OC}_t^*{=}1$ if agent declared
$\kappa{=}\kH$ and trace ended incorrectly.

\textbf{Phase~A---Decision
Transformer~\citep{chen2021decision}.}
Traces $\to$ return-conditioned sequences
$[R_0, s_0, \mathbf{a}_0, \ldots]$, where
$R_t{=}\sum_{t'\geq t}r_{t'}$ and terminal
reward $r_T{=}\text{F1}(\hat{y},y^*)
{-}0.4{\cdot}\text{ECE}$. Training:
$\pi_\theta$ predicts $\mathbf{a}^*_t \mid R_t, s_t$.
Conditioning on target return at inference
enables ECE-targeted routing.

\textbf{Phase~B---GRPO~\citep{guo2025deepseek}.}
$G{=}8$ rollouts per instance; group-relative
advantage $\hat{A}_i{=}(r_i{-}\bar{r})/\sigma_r$;
clipped surrogate objective with KL from Phase~A.
Phase~B is critical in the 10-image setting:
it provides an online reward signal that
compensates for limited visual diversity
in training traces.

\textbf{Reward function:}
\begin{align}
  r &= \text{F1} - 0.4\,\text{ECE}
  + 0.3\,\Delta\text{F1}_\text{BT} \notag\\
  &- 0.05\,(L/L_\text{max})
  + 0.2\,(1{-}|\epsilon_T{-}\epsilon^*_T|). \notag
\end{align}

\textbf{Full objective:}
\begin{equation*}\textstyle
\mathcal{L}{=}\mathcal{L}_\text{rt}
{+}0.4\mathcal{L}_\text{cal}{+}0.2\mathcal{L}_\text{con}
{+}0.2\mathcal{L}_\text{bel}{+}0.3\mathcal{L}_\text{OC}.
\end{equation*}

\textbf{Config:} AdamW lr=$10^{-4}$, cosine decay,
warmup $500$ steps, batch $8{\times}4$; Phase~A
$50$ epochs early stopping (val ECE, patience $5$);
Phase~B $10$ epochs lr=$5{\times}10^{-5}$,
$\beta_\text{KL}{=}0.04$. Single A100 40GB,
${\approx}8$--$10$h. \plantswarm absent at inference.

\subsection{Active Learning for Cross-Crop Convergence}

$\epsilon_t$---not $\alpha_t$---drives
expert queries: $\epsilon_t$ indexes
reducible uncertainty (expert labels provide
useful signal); $\alpha_t$ indexes irreducible
ambiguity (expert labels provide no routing
improvement and waste effort).

\begin{table}[t]\centering\small
\begin{tabular}{lrcc}
\toprule
\textbf{Iter.} & \textbf{Labels}
& \textbf{Query rate} & \textbf{T3 F1} \\
\midrule
0 (zero-shot) & 0         & 60--65\% & 35--50\% \\
1             & ${\approx}500$   & 35--40\% & 55--65\% \\
2             & ${\approx}800$   & 15--20\% & 70--78\% \\
3 (converged) & ${\approx}950$   & ${<}8\%$ & \textbf{80--87\%} \\
\midrule
Supervised & 5,000--10,000 & & \\
\bottomrule
\end{tabular}
\caption{Cross-crop active learning convergence.
  The geo prior $\hat{P}(d\mid s)$ from
  \texttt{db.symptoms.state\_counts} reduces
  zero-shot query rate by providing
  pre-calibrated routing before any
  crop-specific labeled example is seen.}
\label{tab:active}
\end{table}

%% ---------------------------------------------------------------
\section{Experiments}
\label{sec:exp}

\subsection{Headline: seed vs.\ enhanced \pathomedb}
\label{sec:headline}

The primary experiment is a paired comparison
of two \observe checkpoints trained on
\emph{the same} $101{,}640$ \plantswarm traces,
differing only in which \pathomedb the agents
read from at training time:
\begin{itemize}
  \item \texttt{seed}: \pathomedb v2 with
    Claude-headless-seeded visual blocks and
    auto-derived state/AEZ counts. No
    \texttt{swarm\_observations} populated.
  \item \texttt{enhanced}: same DB plus
    per-class \texttt{SwarmObservations}
    from Phase~3 trace mining.
\end{itemize}
Held-out evaluation uses the full \pv (with
seen/unseen slice) and full \pw, run via
\texttt{evaluate\_pathome.py}. The deltas
reported in Table~\ref{tab:beforeafter} are
emitted directly into the paper by
\texttt{compare\_pathome\_versions.py} as
\texttt{$\backslash$PathomeDelta*} macros, so
re-running the pipeline rewrites the table on
\texttt{$\backslash$input}. Numbers below are
the prediction targets (P5--P7); the macros
will fill in observed values.

\begin{table}[t]\centering\small
\setlength{\tabcolsep}{4pt}
\begin{tabular}{lccc}
\toprule
\textbf{Metric} & \textbf{Seed} & \textbf{Enhanced} & $\boldsymbol{\Delta}$ \\
\midrule
T3 macro F1 (full \pv)        & --- & --- & $\geq{+}0.03$ \\
T3 F1 (seen)                  & --- & --- & $\geq{+}0.03$ \\
T3 F1 (unseen)                & --- & --- & $\geq{+}0.02$ \\
ECE (full \pw)                & --- & --- & $\leq{-}0.02$ \\
Avg.\ path length (held-out)  & --- & --- & $\leq{-}0.5$ \\
Backtrack rate                & --- & --- & $\leq{-}0.05$ \\
High-confidence rate          & --- & --- & $\geq{+}0.04$ \\
\bottomrule
\end{tabular}
\caption{Headline before-vs-after table. Seed
  is the Claude-only \pathomedb; Enhanced
  adds \texttt{SwarmObservations} from
  $101{,}640$ trace aggregation. Targets are
  the falsifiable predictions P5--P7 reformulated
  for the seed$\to$enhanced design; observed
  values fill in via
  \texttt{$\backslash$input{auto\_pathome\_metrics}}.}
\label{tab:beforeafter}
\end{table}

\subsection{PlantSwarm Ablation}

\begin{table}[t]\centering\small
\setlength{\tabcolsep}{3pt}
\begin{tabular}{lcccc}
\toprule
\textbf{Variant} & \textbf{Ctx}
& \textbf{Gate} & \textbf{BT}
& \textbf{Purpose} \\
\midrule
Fixed Chain          & \checkmark & $\times$ & $\times$ & Lower bound \\
FC + Full Ctx        & \checkmark & $\times$ & $\times$ & Context value \\
Free, No Gate        & \checkmark & $\times$ & \checkmark & Gate value \\
Free, No Backtrack   & \checkmark & \checkmark & $\times$ & BT value \\
3-Agent Swarm        & \checkmark & \checkmark & \checkmark & Agent count \\
\textbf{\plantswarm} & \checkmark & \checkmark & \checkmark & Full system \\
\bottomrule
\end{tabular}
\caption{Factorial ablation over all $101{,}640$ traces.
  Adjacent variants differ by one component.}
\label{tab:ablation}
\end{table}

\subsection{PlantSwarm Main Results}

\begin{table*}[t]\centering\scriptsize
\setlength{\tabcolsep}{2pt}
\begin{tabular}{lcccccc}
\toprule
\multirow{2}{*}{\textbf{Method}} &
\multicolumn{2}{c}{\textbf{T3 F1}} &
\multicolumn{2}{c}{\textbf{T2 F1}} &
\multirow{2}{*}{\textbf{ECE$\downarrow$}} &
\multirow{2}{*}{\textbf{TPCP$\downarrow$}} \\
\cmidrule(lr){2-3}\cmidrule(lr){4-5}
& \pv (Easy) & \pw (Hard)
& \pv (Easy) & \pw (Hard) & & \\
\midrule
Single VLM zero-shot      & --- & --- & --- & --- & --- & --- \\
DeepSeek-R1~\citep{guo2025deepseek}
                          & --- & --- & --- & --- & ${\approx}0.31$ & --- \\
MVPDR~\citep{tqwei2024mvpdr}
                          & --- & --- & --- & --- & --- & --- \\
Snap \& Diagnose~\citep{tqwei2024snapdiagnose}
                          & --- & --- & --- & --- & --- & --- \\
SAGE + KB ($k{=}16$)~\citep{sage2025}
                          & --- & --- & --- & --- & --- & --- \\
Fixed Chain     & --- & --- & --- & --- & --- & --- \\
Fixed Chain+Ctx& --- & --- & --- & --- & --- & --- \\
\textbf{\plantswarm} & \textbf{---} & \textbf{---}
                          & \textbf{---} & \textbf{---}
                          & ${\approx}\mathbf{0.19}$ & \textbf{---} \\
\bottomrule
\end{tabular}
\caption{\plantswarm evaluated on full \pv (easy) and full \pw (hard).
  \pv and \pw are entirely held out: no training images from either.
  TPCP $= \bar\Omega N_\text{img}/N_\text{correct}$.}
\label{tab:plantswarm-main}
\end{table*}

\subsection{OBSERVE Main Results}

\begin{table*}[t]\centering
\setlength{\tabcolsep}{3pt}
\resizebox{\textwidth}{!}{%
\begin{tabular}{lccccc}
\toprule
\textbf{Method} &
\textbf{ECE\,\pv$\downarrow$} &
\textbf{ECE\,\pw$\downarrow$} &
\textbf{T3\,\pv} &
\textbf{T3\,\pw} &
\textbf{T3$^\dagger$} \\
\midrule
Single VLM       & ${\approx}0.35$& ${\approx}0.41$& ---& ---& --- \\
SAGE + KB        & ${\approx}0.19$& ${\approx}0.38$& ---& ---& --- \\
DeepSeek-R1      & ${\approx}0.22$& ${\approx}0.31$& ---& ---& --- \\
\plantswarm      & ${\approx}0.19$& ${\approx}0.30$& ---& ---& --- \\
+ temp.\ scaling & ${\approx}0.16$& ${\approx}0.27$& ---& ---& --- \\
\observe~(no OC)         & ${\approx}0.13$& ${\approx}0.20$& ---& ---& --- \\
\observe~(no refs)       & ${\approx}0.12$& ${\approx}0.21$& ---& ---& --- \\
\observe~(no geo)        & ${\approx}0.13$& ${\approx}0.23$& ---& ---& --- \\
\observe~(no \pathomedb) & ${\approx}0.16$& ${\approx}0.26$& ---& ---& --- \\
\textbf{\observe~(full)} & ${\approx}\mathbf{0.12}$& ${\approx}\mathbf{0.17}$&
                  \textbf{---}& \textbf{---}& ${\approx}\mathbf{0.43}$ \\
\bottomrule
\end{tabular}
\quad
\begin{tabular}{lccc}
\toprule
\textbf{Method} &
\textbf{OC F1$\uparrow$} &
\textbf{Esc F1$\uparrow$} &
\textbf{Tok$\downarrow$} \\
\midrule
Single VLM       & --- & --- & --- \\
SAGE + KB        & --- & --- & --- \\
DeepSeek-R1      & --- & --- & ${\approx}8{,}000$ \\
\plantswarm      & --- & --- & ${\approx}4{,}200$ \\
+ temp.\ scaling & --- & --- & ${\approx}4{,}200$ \\
\observe~(no OC)         & ---            & --- & ${\approx}700$ \\
\observe~(no refs)       & ${\approx}0.62$& --- & ${\approx}700$ \\
\observe~(no geo)        & ---            & --- & ${\approx}700$ \\
\observe~(no \pathomedb) & ---            & --- & ${\approx}700$ \\
\textbf{\observe~(full)} & ${\approx}\mathbf{0.82}$&
                  ${\approx}\mathbf{0.85}$& ${\approx}\mathbf{720}$ \\
\bottomrule
\end{tabular}}%
\caption{\observe on full \pv and \pw.
  Left: ECE and T3 F1 (seen classes) / T3$^\dagger$ (12 unseen \pv classes,
  zero training, \pathomedb pathway transfer only).
  Right: overconfidence detection, escalation, and token efficiency.
  Full \observe surpasses \plantswarm teacher on both benchmarks.}
\label{tab:observe-main}
\end{table*}

\paragraph{Ablation interpretation.}
The \textit{no refs} row isolates the
contribution of \texttt{ReferenceLibrary} retrieval:
geo-weighted image-to-image comparison drives
overconfidence detection quality, and text-based
self-confidence is substantially weaker.
The \textit{no \pathomedb} row removes
\texttt{SymptomLibrary} access entirely, leaving
\observe with image+context only; the resulting
ECE degradation isolates the joint contribution
of seed visual + state-prior + reference-retrieval.

\paragraph{Coverage of the unseen \pv slice.}
The unseen-class slice is now defined dynamically
as $\pv\,\setminus\,\bugwood$ post-normalisation
rather than fixed at $12$. We report the cardinality
of this slice in the \texttt{seen\_unseen} field
of \texttt{evaluate\_pathome.py}'s output JSON
alongside the per-slice T3 F1.

\subsection{Epistemic/Aleatoric Decomposition}

\begin{table}[t]\centering\small
\begin{tabular}{lc}
\toprule
\textbf{Validation} & \textbf{Expected} \\
\midrule
$\rho(\epsilon_t, \text{BT resolved})$ & $+0.45$--$+0.55$ \\
$\rho(\alpha_t, \text{expert disagr.})$ & $+0.40$--$+0.52$ \\
$\Delta$F1, high-$\epsilon_t$ backtrack & $+7$--$+11$ \\
$\Delta$F1, high-$\alpha_t$ backtrack & ${<}+2$ \\
\midrule
Layer-3 geo prior path length reduction & ${\approx}21\%$ \\
\bottomrule
\end{tabular}
\caption{Uncertainty decomposition and geospatial
  routing efficiency. The near-zero gain from
  high-$\alpha_t$ backtracks confirms aleatoric
  uncertainty is genuinely irreducible, validating
  the active learning query strategy.}
\label{tab:decomp}
\end{table}

%% ---------------------------------------------------------------
\section{Discussion}
\label{sec:discussion}

\paragraph{Inverting the paradigm.}
Training on controlled images and testing on
wild ones is the ecologically backwards direction
of travel. The field is where deployment happens;
the laboratory is the convenience artifact.
Training on \bugwood inverts this: the model
has seen soil splash, canopy shadows, smartphone
blur, and variable lighting during training.
PlantVillage's uniform white backgrounds and
single-leaf framing are \emph{easier} than
the training distribution, not harder, which
explains why ECE on PlantVillage ($0.12$) is
lower than on PlantWild ($0.17$) despite both
being entirely unseen.

\paragraph{One field image per $280$ test images.}
The $1{:}280$ training-to-test ratio is not
a limitation---it is the point. Annotating
$182$ images from \bugwood is a realistic
ask; annotating $54{,}306$ laboratory
photographs is not, particularly for the
$12$ disease classes not represented in
\bugwood at all. \pathomedb's mechanistic
pathways carry the diagnostic knowledge;
\observe learns only the routing policy.
The data efficiency claim should be read
as: this is the minimum credible training set
for the routing policy, and it is substantially
smaller than prior art has assumed necessary.

\paragraph{The seed--enhance--evaluate loop.}
\pathomedb's contribution is structural, not
encyclopaedic. We do not claim that the
$484$ Claude-seeded visual blocks are better
than what a domain expert would write; we
claim that they are good \emph{enough} to bootstrap
trace generation, and that the trace-derived
\texttt{SwarmObservations} measurably improve
\observe at downstream evaluation. The
seed$\to$enhanced delta is the contribution.
A negative or null delta (P5--P7 falsification)
would indicate either that the swarm is too
unreliable to inform the KB, or that
\observe is insensitive to KB content beyond
what's already in the seed---both
falsifiable failure modes by design.

\paragraph{Cross-crop zero-shot transfer.}
T3 F1 on the unseen-\pv slice (diseases
\observe has never encountered in training)
remains a sharp test, but the slice is now
defined dynamically by class-name intersection
between \bugwood usable CSV and \pv after
normalisation. The exact cardinality is reported
in the per-run output JSON. Transfer no longer
relies on a hand-coded mechanistic-pathway
layer; it relies on the visual block (Claude
seed) for the unseen diseases plus the
\texttt{confusion\_diseases} field for
adjacent training classes.

\paragraph{Limitations.}
\label{sec:limits}
The \bugwood IPMNet CSV is US-only ($47/49$
states represented after the $\geq{10}$-row
filter). Layer-3-style monthly priors are
unsupported (no capture date in the export);
the geo prior reduces to $\hat{P}(d\mid s)$
at state granularity. The coarse-fallback FAO
AEZ mapping resolves the entire US footprint
into two zones (TMP, STM); a real GAEZ
shapefile recovers $\sim$$17$ zones via
\texttt{PATHOME\_AEZ\_SHAPEFILE}, but the
EPPO Pearson-$r$ validation of the previous
formulation is not reproducible at this geo
resolution and is dropped. About half of
the $484$ admitted classes appear in only one
state in the CSV (no spatial-variance signal);
the geo prior's contribution concentrates on
the multi-state classes. State-of-the-data
distribution is heavily skewed (NC $33\%$,
ME $18\%$, KY $7\%$), so per-state priors
inherit a survey bias that should be
acknowledged in any deployment outside the
top-data states.

%% ---------------------------------------------------------------
\section{Conclusion}
\label{sec:conclusion}

We have argued and demonstrated that training
on geo-tagged field images is strictly preferable
to training on controlled laboratory images when
the deployment target is field conditions, and
that the curated portion of the knowledge base
need not be assembled by hand. \textbf{\pathomedb v2}
is a two-store, symptom-centric, geo-aware
knowledge base whose visual content is seeded by
\texttt{claude -p} and whose behavioural content
is mined from $101{,}640$ \plantswarm routing
traces. \textbf{\plantswarm} produces richer
behavioural demonstrations from field images
than from controlled photographs (the
self-declared/behavioural-precision gap widens
on field ambiguity), generates the per-class
state-distribution data that grounds the geo
prior empirically, and does so from a corpus
that requires no manual curation beyond
filtering the IPMNet CSV. \textbf{\observe}
trains twice---once against the seed-only
\pathomedb, once against the swarm-enhanced
\pathomedb---on the same trace set, and we
report the seed$\to$enhanced delta on full
\pv (with seen/unseen slice) and full \pw.

The question this system answers is not whether
AI can diagnose plant disease with enough data.
The question is whether a multi-agent VLM swarm,
fed back to the knowledge base it consulted, can
make the next round of agents measurably more
efficient---and how much editorial seeding
is actually needed to get the loop started.

\section*{Ethics Statement}

\bugwood IPMNet images are publicly available
under academic and extension-service terms.
\pathomedb retains location at US-state
granularity (the resolution available in the
CSV export); no per-image GPS that could
identify an individual farm or farmer is
ever stored. The Claude-headless seed is a
\textit{descriptive} resource (``what does
\textit{Botrytis} on strawberry look like?'')
and is not deployed as a treatment
recommender; production deployment requires
trained extension workers with local
regulatory knowledge. We acknowledge that
the IPMNet CSV is US-only and survey-biased
(top three states $=\,57\%$ of records) and
release the per-state coverage statistics
alongside \pathomedb so users can assess
reliability before deployment outside the
top-data states. Per-class T3 F1 and ECE
are released alongside the aggregate metrics.

\section*{Limitations}

The \bugwood IPMNet CSV is US-only at state
granularity; international deployment requires
either a different export with finer GPS or
a separate regional knowledge base.
The coarse-fallback FAO AEZ mapping resolves
the US footprint into two zones; the full
$\sim$$17$-zone resolution requires
\texttt{PATHOME\_AEZ\_SHAPEFILE} pointing at
a real GAEZ shapefile.
About half of the $484$ admitted classes
appear in only one state, contributing no
spatial-variance signal; the geo prior is
informative on the multi-state subset only.
\texttt{ReferenceLibrary} contains $1{,}452$
images---larger than the previous
formulation's $78$, but rare disease
stages may still retrieve suboptimal
references for visually-confounded classes.
\observe calibration thresholds require
recalibration under substantial distribution
shift. $\kappa$-circularity partially
propagates label noise; GRPO Phase~B
mitigates but does not eliminate it. The
Claude-headless seed quality varies with
disease prevalence in the model's training
data; rare or recently-described diseases
are more likely to return empty visual
fields.

\section*{Acknowledgments}
[Omitted for anonymous review.]

\bibliographystyle{plainnat}
\bibliography{pathome3}

%% ---------------------------------------------------------------
\appendix

\section{Bugwood CSV Filtering}
\label{app:selection}

The IPMNet export
(\texttt{BugWood\_Diseases.csv}, $19{,}749$
rows) is normalised by
\texttt{scripts/filter\_bugwood\_csv.py}:
\begin{enumerate}
  \item drop rows with empty \texttt{Host Name}
    ($2{,}171$);
  \item map \texttt{Host Name} via
    \texttt{BUGWOOD\_EXACT\_CROP\_MAP} (60+
    crop synonyms; non-host taxa dropped);
  \item strip parenthetical Latin from
    \texttt{Subject Display Name}; title-case;
  \item drop rows with empty/unrecognised
    \texttt{Location};
  \item group by $(\text{NormCrop},\,\text{NormDisease})$
    and keep classes with $\geq{10}$ rows.
\end{enumerate}
Output: \texttt{BugWood\_Diseases\_usable.csv},
$11{,}513$ rows over $\mathbf{484}$ classes.
Five derived columns are appended
(\texttt{NormCrop}, \texttt{NormDisease},
\texttt{StateLat}, \texttt{StateLon},
\texttt{AezCode}). The first $10$ rows per
class (deterministic ordering by
\texttt{Image Number}) are admitted by
\texttt{BugwoodLoader}; $7$ feed \plantswarm
trace generation, $3$ become \texttt{ReferenceLibrary}
entries.

\texttt{select\_diverse\_subset} (CLIP
$k$-medoids over the candidate pool) is still
exposed for callers that want to apply
diversity sampling \emph{within} a class's
post-filter pool, but is not used by default;
the deterministic Image-Number ordering is
sufficient for reproducibility.

\section{Geo prior at state granularity}
\label{app:layer3}

State centroids are population-weighted (US
Census 2020) for the $50$ states + DC + US
territories that occasionally appear in IPMNet.
The reverse-lookup in
\texttt{pathome.database.\_nearest\_state}
resolves $(\lambda,\,\phi)$ to the closest
centroid in $\ell_2$. The prior reads
$\hat{P}(d\mid s) =
\mathrm{count}(d, s) / \sum_{d'}\mathrm{count}(d', s)$
directly off the symptom profile's
\texttt{state\_counts}. State cells with
$<3$ observations are marked \texttt{sparse};
\observe queries the global pathogen prior
in those cases. The previous AEZ-month grid
and EPPO Pearson-$r$ validation are dropped
because the IPMNet export lacks the GPS and
date fields required to construct them.

\section{Training Configuration}
\label{app:training}

\textbf{Phase A (DT):} AdamW $\text{lr}{=}10^{-4}$,
cosine decay, warmup $500$ steps, batch $32$
($8{\times}4$ accumulation), $50$ epochs,
early stopping patience $5$ (validation ECE).
Train/val/held-out split: 80/10/10 of the
$101{,}640$ \plantswarm traces.
$1\times$ A100 80GB, ${\approx}10$--$14$h
per checkpoint (the seed-DB and enhanced-DB
runs are sequential in
\texttt{submit\_pathome\_phase4\_train.sh}).

\textbf{Phase B (GRPO):} AdamW $\text{lr}{=}5{\times}10^{-5}$,
$G{=}8$ rollouts, $10$ epochs, $\beta_\text{KL}{=}0.04$,
clip $\epsilon{=}0.2$, reference policy frozen from Phase~A.
$1\times$ A100 80GB, ${\approx}6$--$8$h per checkpoint.

\section{Falsifiable Predictions}
\label{app:predictions}

\begin{table}[H]\centering\small
\begin{tabular}{lp{3.4cm}c}
\toprule
\textbf{P} & \textbf{Quantity} & \textbf{Target} \\
\midrule
P1 & $\rho(L, H[\hat{y}])$ & $+0.42$--$+0.55$ \\
P2 & $\Delta$Acc, PathogenAgent post-BT & $+9$ F1 \\
P3 & acc(early) $-$ acc(extended) & $+12$ F1 \\
P4 & \observe ECE (\pv / \pw) & ${\leq}0.15/0.20$ \\
P5 & Enhanced$-$seed $\Delta$T3 macro F1 (full \pv) & ${\geq}{+}0.03$ \\
P6 & Enhanced$-$seed $\Delta$ECE (full \pw) & ${\leq}{-}0.02$ \\
P7 & Enhanced$-$seed $\Delta$ avg.\ path length & ${\leq}{-}0.5$ \\
\bottomrule
\end{tabular}
\caption{Seven falsifiable predictions. P5--P7
  reformulated for the seed$\to$enhanced design;
  observed values fill in via
  \texttt{compare\_pathome\_versions.py} \texttt{$\backslash$PathomeDelta*} macros.}
\label{tab:predictions}
\end{table}

\end{document}
